-	Quellenarbeit gut, auch inhaltlich
-> wurde noch ein 2. Mal erwähnt :D
-	Hat gefehlt: Forschungsschwerpunkt:
o	Wie kann man Diskriminierung technisch umgehen
-> Thema explanable AI
o	Aktuelle Rechliche Themen, z.B: AI-Act
-> für Hochrisikosysteme
o	Mehr Beispiele (z.B. in 3.3.2)
	Rückfälle Straftäter USA
	Amazon KI
	An den relevanten  Stellen hinzufügen, in den anderen eher weniger
-> z.B. wo das Thema Diskriminierung Probleme mit sich bringt
	Mir ist am Ende der Platz ausgegangen :/
o	Umgangssprache
	Einzelne Aspekte, nicht generell
	S. 8, K. 3.3.1 – „sich schmerzlich erinnern können“
-> negative Erinnerungen / Eindrücke in Bezug darauf
	10-tausende Datensätze
-> große Mengen
o	Graphiken / Schaubild
	Z.B. Kapitel Deep Learning
-> War zuviel Fließtext
	Abbildung zum auflockern
o	Kapitel 2.2 – 2.2.1 / 2.2.2
-> Text vor den Kapiteln
o	Mehr Quellen angeben
-> „in einem Szenario“, welches?
-> war eigentlich auf das vorhergehende bezogen? – ich hatte ein Szenario aufgestellt und das für mehrere Kapitel genommen
